\section*{Capítulo \ref{ch:b3:pericyclic} --- Reações pericíclicas}

\begin{solution}{exo:b3:pericyclic:1}
Diels--Alder: \omterm{def:b3:pericyclic:families}{cicloadição} [4+2]. Abertura do ciclobuteno: eletrocíclica. Cope: sigmatrópico
[3,3]. Perda de $\ce{SO2}$ pelo sulfoleno: queletrópica. Ozônio e alqueno: \omterm{def:b3:pericyclic:dipolar}{cicloadição 1,3-dipolar}
([3+2]). Deslocamento $[1,5]$-H: sigmatrópico.
\end{solution}

\begin{solution}{exo:b3:pericyclic:2}
Hexatrieno (6 elétrons): térmico \omterm{def:b3:pericyclic:rotation-modes}{disrotatório}, fotoquímico \omterm{def:b3:pericyclic:rotation-modes}{conrotatório}.
Butadieno (4): térmico \omterm{def:b3:pericyclic:rotation-modes}{conrotatório}, fotoquímico \omterm{def:b3:pericyclic:rotation-modes}{disrotatório}.
\end{solution}

\begin{solution}{exo:b3:pericyclic:3}
Quatro elétrons, térmica: abertura \omterm{def:b3:pericyclic:rotation-modes}{conrotatória}. Os dois grupos metila, em faces opostas
do anel, giram ambos para fora: (\textit{E},\textit{E})-hexa-2,4-dieno.
\end{solution}

\begin{solution}{exo:b3:pericyclic:4}
A reação [2+2] \omterm{def:b3:pericyclic:faciality}{suprafacial} tem $4q$ elétrons: no estado fundamental, a HOMO de um
eteno e a LUMO do outro se sobrepõem com uma extremidade ligante e uma antiligante.
No estado excitado, o $\pi^*$ monoocupado de uma molécula faz o papel da
HOMO, e as duas extremidades se correspondem.
\end{solution}

\begin{solution}{exo:b3:pericyclic:5}
O \omterm{def:b3:point-groups:mirror}{plano especular} se conserva. Hexatrieno: $\psi_1$ S, $\psi_2$ A, $\psi_3$ S ocupados; $\psi_4$ A, $\psi_5$ S, $\psi_6$ A vazios.
Ciclo-hexadieno, em ordem: $\sigma$ S, $\pi_1$ S, $\pi_2$ A ocupados; $\pi_3^*$ S, $\pi_4^*$ A, $\sigma^*$ A vazios. Ligando por
simetria em ordem, $\psi_1 \to \sigma$, $\psi_2 \to \pi_2$, $\psi_3 \to \pi_1$: de ocupado para ocupado, o fechamento \omterm{def:b3:pericyclic:rotation-modes}{disrotatório}
é termicamente permitido.
\end{solution}

\begin{solution}{exo:b3:pericyclic:6}
Seis elétrons sob luz: \omterm{def:b3:pericyclic:rotation-modes}{conrotatório}, de modo que os dois grupos metila terminais terminam em
faces opostas: \textit{trans}-5,6-dimetilciclo-hexa-1,3-dieno (a reação térmica dá
o isômero \textit{cis}).
\end{solution}

\begin{solution}{exo:b3:pericyclic:7}
O hidrogênio de C5 passa para C1 do dieno, seu vizinho ao redor do anel,
por um estado de transição de seis elétrons ($\sigma2_s + \pi4_s$), de modo \omterm{def:b3:pericyclic:faciality}{suprafacial}: uma migração
dá o 1-metil-, uma segunda o 2-metilciclopentadieno. Uma cadeia de cinco carbonos permite que o
hidrogênio permaneça em uma face, o que o anel pequeno impõe de qualquer modo.
\end{solution}

\begin{solution}{exo:b3:pericyclic:8}
$\pi2_s$ (alqueno) $+ \pi2_a$ (C=C da cetena, atacada em faces opostas através do
$\pi^*$ C=O perpendicular): o número de componentes $(4q + 2)_s$ é um (o alqueno), ímpar:
termicamente permitida.
\end{solution}

\begin{solution}{exo:b3:pericyclic:9}
Na cadeira, a C=C do grupo vinila (átomos 1--2), o oxigênio (3), o carbono
$\ce{OCH2}$ (4) e a C=C alílica (5--6); a ligação O--C se rompe e a ligação C1--C6 se forma:
o produto é o pent-4-enal,
$\ce{OHC-CH2-CH2-CH=CH2}$.
\end{solution}

\begin{solution}{exo:b3:pericyclic:10}
Oito elétrons correspondem a $4q$; com um componente \omterm{def:b3:pericyclic:faciality}{antarafacial}, o arranjo tem topologia de
Möbius, aromática com $4q$ elétrons: o estado de transição é aromático e a
reação é termicamente permitida (como o fechamento \omterm{def:b3:pericyclic:rotation-modes}{conrotatório} do octatetraeno).
\end{solution}

\begin{solution}{exo:b3:pericyclic:11}
Emparelhar os maiores coeficientes une N3 ao carbono terminal: o
triazol 1,4-dissubstituído. Sem catalisador, os coeficientes diferem pouco e
formam-se os dois isômeros, 1,4 e 1,5; o cobre(I) torna a reação gradual, por meio de um
acetileto de cobre, e dá apenas o isômero 1,4.
\end{solution}

\begin{solution}{exo:b3:pericyclic:12}
O hidrogênio faz ponte entre as extremidades da SOMO de um radical de $j$ átomos, $k = (j + 1)/2$. Seus coeficientes
terminais são proporcionais a $\sin(k\pi/(j + 1))$ e $(-1)^{k+1}\sin(k\pi/(j + 1))$. Para $j = 7$, $k = 4$: sinais opostos, logo
os lobos de mesmo sinal ficam em faces opostas nas duas extremidades: \omterm{def:b3:pericyclic:faciality}{antarafacial}. Para
$j = 5$, $k = 3$: o mesmo sinal na mesma face: \omterm{def:b3:pericyclic:faciality}{suprafacial}.
\end{solution}

\begin{solution}{pb:b3:pericyclic:1}
\textbf{1.} A ligação $\sigma$ C9--C10 do anel B; com o 5,7-dieno, ela dá um hexatrieno,
C10=C5--C6=C7--C8=C9.

\textbf{2.} Seis elétrons (as duas ligações $\pi$ e a ligação $\sigma$): uma abertura de anel eletrocíclica.

\textbf{3.} Termicamente, um ciclo-hexadieno e seu hexatrieno aberto só se equilibram
lentamente, por uma barreira alta, e o anel é o lado mais estável; o
dieno absorve luz ultravioleta, e seu estado excitado se abre em
picossegundos.

\textbf{4.} \omterm{def:b3:pericyclic:rotation-modes}{Conrotatória}: no estado excitado, a LUMO monoocupada do sistema
trieno ($k = 4$) tem extremidades de sinais opostos.

\textbf{5.} A ligação C6=C7 vem do anel: suas duas continuações da cadeia, C5 e C8,
estavam do mesmo lado dela, e assim permanecem: \textit{Z}.

\textbf{6.} A abertura térmica, permitida apenas de modo \omterm{def:b3:pericyclic:rotation-modes}{disrotatório}, tem uma barreira muito além do que
o calor do corpo pode fornecer, e seria ascendente, em direção ao trieno menos estável.

\textbf{7.} De C19 (a metila em C10) para C9: um deslocamento sigmatrópico $[1,7]$-H.

\textbf{8.} Oito: os seis elétrons $\pi$ do trieno e os dois da ligação C--H.

\textbf{9.} $\sigma2 + \pi6$. De modo \omterm{def:b3:pericyclic:faciality}{suprafacial}, ambos seriam s, com um componente $(4q + 2)_s$ de cada lado:
par, proibida. Com o trieno \omterm{def:b3:pericyclic:faciality}{antarafacial}, $\sigma2_s + \pi6_a$: um componente contado, ímpar,
permitida.

\textbf{10.} Sete átomos em uma cadeia helicoidal de configuração \textit{Z} permitem que o hidrogênio passe de uma face para
a outra; uma cadeia de três átomos não consegue se torcer tanto.

\textbf{11.} Möbius (um componente \omterm{def:b3:pericyclic:faciality}{antarafacial}) com oito elétrons, $4q$: aromático.

\textbf{12.} Ele é termicamente permitido; sua barreira é transposta lentamente à temperatura do corpo.

\textbf{13.} Seis elétrons sob luz: de novo \omterm{def:b3:pericyclic:rotation-modes}{conrotatório}, mas ela pode girar as extremidades no
outro sentido e pôr H9 e a metila C19 nas outras faces: um estereoisômero de anel
fechado, o lumisterol.

\textbf{14.} Com a ligação central \textit{E}, C19 e C9 ficam em lados opostos da cadeia e longe
um do outro: o hidrogênio não consegue alcançar.

\textbf{15.} A própria pré-vitamina D$_3$ absorve luz e é convertida em lumisterol e taquisterol,
que não dão vitamina D; atinge-se uma mistura fotoestacionária, que limita a
produção.

\textbf{16.} A abertura ocorre no estado excitado, em femtossegundos a picossegundos;
o deslocamento é uma reação térmica com uma barreira de \qty{85}{kJ/mol}.

\textbf{17.} $t_{1/2} = \ln 2/k = 0.693/\qty{0.023}{h^{-1}} = \qty{30}{h}$.

\textbf{18.} $1 - \eu^{-0.023 \times 24} = 0.42$.

\textbf{19.} $\ln 10/k = 2.303/\qty{0.023}{h^{-1}} = \qty{100}{h}$.

\textbf{20.} $k(\qty{293.15}{K}) = k(\qty{310.15}{K})\,\eu^{-(E_a/R)(1/293.15 - 1/310.15)} = 0.023 \times \eu^{-1.91} = \qty{3.4e-3}{h^{-1}}$.

\textbf{21.} $2.303/\qty{3.4e-3}{h^{-1}} = \qty{680}{h}$, cerca de 28 dias.

\textbf{22.} A etapa térmica é sete vezes mais rápida a \qty{37}{\celsius} do que a \qty{20}{\celsius}, de modo que a
pré-vitamina formada em uma manhã ao sol é convertida em poucos dias.

\textbf{23.} O esqueleto esteroide, com seus numerosos estereocentros, vem pronto do
esterol; a luz e o calor realizam então as duas etapas pericíclicas exatamente como na
pele.

\textbf{24.} À temperatura do corpo, 90\,\% da pré-vitamina D$_3$ se transforma em vitamina D$_3$ em cerca de
\qty{100}{h}, quatro dias.
\end{solution}
